\section{задание: Суммаризация (Сокращение текста)}
\subsection{исходные данные}
\indent Тема исходного текста: \textbf{Компьютерное зрение} \\
\indent Формат вывода: \textbf{Таблица сравнения} \\
\indent Степень сжатия: \textbf{Средняя} \\
\indent Фокус внимания: \textbf{Скорость и точность} \\
% \indent Ссылка на статью с текстом: \textbf{ \\ \indent \indent https://habr.com/ru/companies/otus/articles/921402/} \\
\indent Ссылка на статью с текстом: \textbf{ \\ \indent \indent https://habr.com/ru/articles/946080/ (1500+ слов)} \\
\newpage
\noindent Фрагмент текста:
\begin{verbatim}
		Инженерные чертежи содержат десятки типов размеров и допусков:
		линейные и угловые, радиальные и диаметральные, справочные и
		базовые, а также геометрические характеристики вроде плоскостности
		или перпендикулярности. Для инженера это привычный язык, но
		для автоматической обработки такие документы представляют
		серьёзную задачу.

		С появлением мультимодальных LLM возник вопрос: способны ли
		они заменить классические OCR-решения и специализированные
		пайплайны постобработки? Насколько хорошо современные модели
		справляются с извлечением размеров и допусков из коробки, без
		дообучения? Какие модели дают максимальную точность, какие
		быстрее, а какие дешевле? И можно ли комбинировать разные
		подходы, чтобы повысить результат?

		Чтобы ответить на эти вопросы, мы протестировали шесть актуальных
		vision-LLM на одном наборе реальных механических чертежей и
		сравнили их точность, скорость и стоимость обработки.

		<< ... >>
\end{verbatim}
\newpage
\subsection{базовый промпт}
\begin{verbatim}
					Сделай краткое содержание этого текста: {ТЕКСТ}
\end{verbatim}
\noindent Самый простой и понятный запрос для сжатия текста.
\subsection{журнал промпт-инжиниринга}
\noindent Промпт v1:
\begin{verbatim}
					Сделай краткое содержание этого текста: ...
\end{verbatim}
\newpage
\noindent Результат (YandexGPT):
\begin{figure}[H]
	\centering
	\includegraphics[width=14cm]{z4/a1.png}
\end{figure}
\noindent Анализ ошибок: модель предоставила ответ в другом формате. \\
Корректировка: добавить уточнение в запрос. \\
\newline
\noindent Промпт v2:
\begin{verbatim}
					Сделай краткое содержание текста со степенью сжатия: Средняя.
					Формат вывода содержимого строго только: Таблица сравнения.
					Фокус внимания в содержимом на: Скорость и точность.
					текст: ...
\end{verbatim}
\newpage
\noindent Результат (YandexGPT):
\begin{figure}[H]
	\centering
	\includegraphics[width=15cm]{z4/b1.png}
\end{figure}
\newpage
\noindent Анализ ошибок: модель отклонилась от заданной структуры. \\
Корректировка: использовать другую модель. \\
\newline
\noindent Результат (DeepSeek):
\begin{figure}[H]
	\centering
	\includegraphics[width=17cm]{z4/c1.png}
\end{figure}
\noindent Анализ ошибок: модель предоставила таблицу сравнения. \\
Корректировка: не требуется. \\
\newpage
\subsection{Сравнение моделей:}
\begin{table}[h!]
	\centering
	\begin{tabular}{|l|p{5cm}|p{5cm}|}
		\hline
		\textbf{Критерий} & \textbf{YandexGPT} & \textbf{DeepSeek} \\
		\hline
		Информационная плотность & 5 & 5 \\
		\hline
		Отсутствие искажений & 5 & 5 \\
		\hline
		Следование инструкции & 3 & 5 \\
		\hline
		\textbf{Итог} & 4.33 & 5 \\
		\hline
	\end{tabular}
\end{table}
Модель DeepSeek снова выдала ответ лучше, так как имеет функцию "подумать"\ что позволяет сгенерировать более валидный результат, а модель YandexGPT не следует инструкции и добавляет лишние разделы.
\subsection{Финальный промпт:}
\begin{verbatim}
					Сделай краткое содержание текста со степенью сжатия: {СТЕПЕНЬ}.
					Формат вывода содержимого строго только: {ФОРМАТ}.
					Фокус внимания в содержимом на: {ФОКУС}.
					текст: {ТЕКСТ}
\end{verbatim}
\subsection{Вывод по заданию:}
\indent В данном задании рекомендуется использовать технику <<Уплотнение смысла>>, но мне показадлсь что для сжатия всей статьи в таблицу сравнения достаточно и обычного запроса.